% s-302-p302-m002; original PDF 302, 303, 304, 305
\Aphorism{s-302-p302-m002}{358}
Wie konnte man Schopenhauer wissenschaftlich nehmen! Das ist ein Fall von Kurzsichtigkeit, der in Deutschland endemisch auftritt. In Frankreich scheint man ganz andere Zugänge zu dem grossen Mystagogen entdeckt zu haben, sei es auch nur, dass sich die kleinen \Foreign{décadents} zur letzten Aufpeitschung ihrer abgewirthschafteten Nerven seiner bemächtigen. Absinth, Schopenhauer, Tolstoj; Florenz und Savonarola; ein Occultismus, den das Rückenmark dem weichen Gehirn suggerirt \Dots{3} Jedenfalls wissen die Pariser, wohin sie sich ihren Schopenhauer zu appliciren haben, während wir als wissenschaftliche Biedermänner nach einer Lehre, einem System fahnden. Damit thun wir Schopenhauer selbst schweres Unrecht, sofern nämlich Einer, der in der Wissenschaft mitzählen will, darauf verzichten muss, in vielen anderen Dingen in Betracht zu kommen. Schopenhauer missverstand sich selbst, und zwar in \Foreign{malam} \Foreign{partem}, wenn er sich für einen theoretischen Denker hielt; auch ihn, wie manchen Anderen, lockte das Licht der „wissenschaftlichen Wahrheit“, das gerade für den Philosophen ein Irrlicht ist. Eine Philosophie, die wahr sein will, wird Moralstatistik oder Sinnesphysiologie oder sonst etwas Exactes, aber sie bleibt nicht Philosophie. Kann man sich Heraklit, oder Spinoza, oder Schopenhauer unter die Alternative „wahr oder unwahr“ gebracht denken? Oder genauer, können wir uns einen einzigen Satz ihrer Systeme denken, der \Sperr{nicht} unwahr, Irrthum, Idiosynkrasie, persönliches Geschmacksurtheil wäre? Dennoch, nein gerade desswegen sind diese Männer gross und unersetzlich, während man den Urheber einer wissenschaftlichen Entdeckung streichen kann, ohne die Entdeckung selbst zu schädigen. — Um einen Einzelfall zu nehmen, so ist uns durch Schopenhauer heute noch einmal (vermuthlich zum letzten Mal) die Möglichkeit des \Sperr{religiösen} Lebens eröffnet worden; natürlich nicht als wissenschaftliche Möglichkeit, wie gelegentlich behauptet wird, sondern in einem anderen, individuellen Sinne. Das theoretische Hinterpförtchen, das Kant und Schopenhauer sich zur positiven Metaphysik gelassen haben (nämlich jener Dualismus, wonach für die Welt als „Ding an sich“ das Gegentheil von dem gelten soll, was für die Welt als „Erscheinung“ gilt), das ist uns zu niedrig: da kann Unsereiner aufrechten Ganges nicht hindurch. Aber was hindert uns, statt dieses wissenschaftlichen Schleichweges den offenen ästhetischen Zugang zu wählen? Die prachtvolle Synthesis kühner und leidenschaftlicher Conceptionen, die übersättigte Lösung indischer Wohlgerüche und Phantasmen, die Schopenhauer an Stelle des europäisch verflachten, moralverdünnten Christenthums darzureichen weiss, lässt uns die Religion noch einmal zum Erregungscentrum und Genussmittel hohen, wenn auch nicht ersten Ranges werden. Wir müssten unsere Vernunft und Moralität an der unrechten Stelle haben, wenn wir etwa der dummen theoretischen Ehrlichkeit zu Liebe auf diese zeitweilige Möglichkeit eigenartiger und buntfarbiger Erlebnisse verzichten wollten; sie geht bald genug vorüber. Eines Tages sind wir alt und können nicht mehr an den Brüsten des Lebens saugen; dann heisst es, das Gesammelte ordnen, stilisiren, aufeinanderbeziehen, in ein System bringen, dann entspinnt sich in unserem Gehirn der theoretische Kampf ums Dasein, der mit dem Siege der leitenden Grundgedanken und mit der Ausrottung der unverträglichen Einzelgedanken endet. Bis dahin steht unser Bewusstsein allen einströmenden Inhalten offen, ohne Unterschied, ohne Prüfung der wissenschaftlichen Legitimation; bis dahin ist das ganze Gebiet des historisch Erlebten unserem Nacherleben preisgegeben, vorausgesetzt dass unsere wählerische Zunge nicht Nein sagt. Wir halten es mit der seelischen Kost wie der gesunde Mensch mit der leiblichen: er isst, was ihm schmeckt, ohne zuvor den chemischen Nährwerth feststellen zu lassen. Ob Religion theoretische „Zellen bildet“, was geht das uns an? wir fragen nur, ob sie uns zusagt oder nicht — schlimmstenfalls verderben wir uns den Magen. Wir haben uns „Gott“ nie von der Wissenschaft verbieten lassen; wir liessen ihn nur stehen, weil er schal geworden war. Der Stimmungstrank, den uns Schopenhauer als Religion kredenzt, ist alles eher als schal; folglich kosten wir ihn, berauschen uns an ihm, und lachen über die Wissenschaft, die uns vor Lebensmittel-Verfälschung warnen will.
\par

